\documentclass{lab1report}
\usepackage{float}     % 支持 [H] 固定浮动体位置
\usepackage{booktabs}  % 支持 \toprule 等三线表命令
\usepackage{needspace} % 需要时将完整代码块移至下一页
\usepackage{multirow}  % 表格中跨行显示主机角色和用途

% 封面信息
\course{网络安全}
\reporttitle{实验报告}
\student{\begin{tabularx}{\linewidth}{@{}>{\centering\arraybackslash}X@{}}
江李阳（202400460042）\\[-0.2ex]\hline
王健一（202400460037）\\[-0.2ex]\hline
邱泽辉（202400460046）\\[-0.2ex]\hline
\end{tabularx}}
\instructor{郭山清}
\college{网络空间安全学院}
\major{密码科学与技术2024级一班}
\completiondate{\today}

% 参考资料写入 references.bib；直接引用使用 \labquote[页码]{文献键}。
% \addbibresource{references.bib}

% 四个任务的起始一级标题居中，其他一级标题沿用左对齐样式。
\newcommand{\labinputexperiment}[1]{%
    \titleformat{\section}[block]{\centering\heiti\zihao{3}}{\thesection}{0.75em}{}%
    \input{#1}%
    \titleformat{\section}{\heiti\zihao{3}}{\thesection}{0.75em}{}%
}

\begin{document}
\makecover

\section{实验目的}

\section{实验原理}
\subsection{互联网域名解析}
互联网域名解析过程由四个关键服务器发挥作用，分别为:
\begin{itemize}
    \item 递归域名服务器：位于用户端与域名服务器之间，负责处理用户端发起的域名解析请求。在接收到用户端发起的域名查询请求后，递归解析服务器会逐层向根域名服务器、顶级域名服务器、权威域名服务器发送请求，直到获得域名对应的IP地址，并将查询结果返回给用户端；
    \item 根服务器：是域名系统中域名解析流程的起点，主要作用是提供顶级域名服务器和权威域名服务器信息；
    \item 顶级域名服务器：负责存储与顶级顶级域名相关的信息及其下所有域名的权威服务器的信息，在域名解析流程中，顶级域名服务器负责提供域名的权威服务器地址，帮助递归解析服务器完成定位；
    \item 权威服务器：负责提供域名的所有信息，包括IP地址、邮件服务器等。
\end{itemize}
\subsection{DNS缓存污染攻击}
为了提高域名解析的性能，减轻权威服务器侧的网络流量负载，递归服务器引入了缓存机制，即递归解析服务器在接收到查询的域名相应后，可将其中的信息存储在本地一段时间，缓存时间由权威服务器指定。
在同一网络环境内其他用户所要查询的域名信息已被缓存，那么递归解析服务器可以直接返回缓存数据，无需再向远程的权威服务器重复发送查询请求。

缓存污染攻击的目的是为了篡改域名在缓存中的资源记录，以进一步实现网络钓鱼、传播恶意软件或收集用户敏感信息等。该攻击能实现的根源在于域名协议报文缺乏消息完整性验证的检查。
攻击者通过发送精心构造的DNS响应报文来欺骗递归解析服务器，使其将错误的解析结果存储在缓存中（缓存保存有时间限制）。当递归解析服务器被污染，缩用使用该递归解析服务器的互联网用户在解析该域名时都会被错误重定向。

上述攻击的关键在于，攻击者发送的伪造DNS响应报文能被递归解析服务器所接受，这在于以下字段核验通过：
\begin{itemize}
    \item IP地址匹配：这一条符合，因为响应的源IP地址为权威服务器地址，目的地址为递归解析服务器地址，均已知；
    \item 端口匹配：响应报文的源端口为DNS服务端口，属于常见端口为53，但目的端口为递归解析服务器发送请求时的随机端口，该值未知，需要遍历猜测。
    \item TXID匹配：响应TXID于请求TXID相等，该值无法提前获知，需要遍历猜测。
    \item 问题区域匹配：通过攻击者自行发起DNS解析服务，从而问题区内容有攻击者掌握，这一条符合。
\end{itemize}
以上所有字段核验通过后还要求攻击者伪造的响应比权威服务器的真实响应先到达递归解析服务器。上述16位源端口和16位TXID待遍历，意味着总的遍历次数为
$$
2^{16} \times 2^{16} = 2^{32}
$$
即需要在权威服务器返回真实响应前做完$2^{32}$次枚举，并将所有伪造响应都发送到递归解析服务器，DNS污染攻击才能成功。
\subsection{DNS反射放大攻击}

反射放大攻击是分布式拒绝服务攻击（Distribute Denial of Service ,DDoS）的一种方式，其攻击根源在于网络中存在的大量的域名递归解析服务器，且这些服务器未限用户的目标范围，
对攻击者公开暴露。攻击者通过借助伪造源地址的方式向大量递归解析服务器发起查询请求，这些递归解析服务器将向受害者主机发送大量回复报文，由于回复报文长度远大于查询寻报文长度，
因此造成了流量放大。

在这种攻击场景中，反射器是互联网中大量的开放的递归解析服务器，放大器是由于DNS响应报文远大于请求报文。



\subsection{DNS重绑定攻击}

DNS重绑定攻击(DNS Rebinding)是一种利用域名解析机制发起的网络攻击方法，主要利用域名解析的应答响应报文，规避浏览器的同源防御策略，
获取受害者内网资源。

DNS重绑定攻击的实用场景未受害者处于安全受控的内网环境，攻击者难以直接从外部网络环境对受害者内网资源进行访问。
具体攻击流程为：
\begin{itemize}
    \item 攻击者引导受害者访问其受控制的恶意域名，恶意域名服务器将该域名指向一个公开可访问的主机地址，并将恶意域名资源记录的有效TTL设置为极短。
    \item 恶意网站中包含JS的恶意代码，但是会受到同源策略的限制。
    \item 当资源记录过期时，攻击者控制的恶意域名服务器将该域名解析为受害者所在的内部网络的主机地址，由于此时浏览器认为整个网络访问过程属于同一个域名（这就绕过了浏览器同源策略，其中同指协议、端口、主机都相同），
    将允许恶意JS代码访问受害者内部网络环境。
\end{itemize}


\labinputexperiment{content/task1.tex}
\labinputexperiment{content/task2.tex}
\labinputexperiment{content/task3.tex}
\labinputexperiment{content/task4.tex}

\input{content/discussion.tex}

\appendix
\input{content/append/key-code.tex}

\end{document}
